% \subsection{Introduction}
% \label{uvm:sec:intro}

Modern Large Language Model (LLM) inference workloads
    are increasingly bottlenecked by GPU memory capacity:
    flagship model sizes grow roughly $410\times$ every two years,
    while GPU memory only doubles in the same period~\cite{ai-memory-wall}.
A single high-end accelerator can no longer hold even a moderately sized model
    together with its KV caches.
To overcome the GPU memory wall,
    modern applications rely on application-driven \textbf{manual memory management} that
    stages weights and activations in host CPU memory~\cite{flexgen, infinigen, headinfer, vllm}.
Each such framework re-implements its own memory manager
    and dictates when and where to move data.

Manual memory management poses two major challenges.
First, offload paths are not overhead-free:
    they demand careful profiling so that copy cost stays hidden behind computation.
Second, GPU kernels and allocator logic must be rewritten for
    \emph{every} application that adopts offloading,
    and tuned carefully to avoid fatal GPU out-of-memory (OOM) errors.
That refactoring is costly: over 7K lines of code in DLRM alone~\cite{fbgemm_tbe_training_docs}.


In contrast, driver-managed automatic memory such as CUDA Unified Virtual Memory (UVM)
    offers another way past the GPU memory wall.
UVM lets the GPU transparently access CPU and peer-GPU memory,
    enabling oversubscription without application-specific offload code---in principle,
    it could retire much of the bespoke machinery that manual paths require.
In practice, however, UVM has been considered too slow for production:
    PyTorch's own documentation advises against \texttt{cudaMallocManaged},
    the default UVM allocator.
Prior work reports that UVM is over $5\times$ slower than
    application-level offloading on LLM workloads~\cite{infinigen}.
The gap is two-fold.
First, frameworks such as PyTorch and vLLM
    offer no first-class UVM integration,
    so managed allocations bypass their caching and pooling optimizations
    and incur memory traffic that an application-aware allocator could elide.
Second, the driver has no visibility into high-level access patterns,
    so page faults and migrations are triggered late and often evict data
    that the application is about to reuse.

We present {\uvmname}, an application-aware UVM substrate for PyTorch
    that closes this performance gap and turns UVM
    into a practical tool for LLM inference.
{\uvmname} is built around three pieces that propagate
    application-level intent down to the UVM driver
    without changing model code.
First, a \emph{custom UVM-aware caching allocator}
    proactively returns cached memory to the driver when the GPU is oversubscribed,
    avoiding the OOM kills that naive UVM caching suffers.
Second, a \emph{discard-based eviction} strategy
    marks partially-free cached blocks as evictable
    so the driver can reclaim them without copying their data back to host.
% Third, \emph{prefetching with location hints}
%     hides GPU page-fault latency on reuse
%     by warming the next layer's tensors before they are needed.

Even with these optimizations, page-fault migration still degenerates into thrashing
    once the working set substantially exceeds GPU memory:
    every fault-in forces an eviction of a page that is itself about to be re-used,
    and discard reduces eviction cost but does not break the eviction--re-fault cycle.
We therefore explore two complementary mechanisms.
\emph{Direct CPU-memory access} leaves cold pages on host
    and lets the GPU read them in place over the interconnect,
    short-circuiting the eviction--re-fault cycle entirely.
\emph{Access-counter--based migration} defers fault-in until a page
    has been touched above a threshold,
    filtering one-shot accesses that would otherwise pay full migration cost
    only to be evicted shortly after.
Together, they reveal that the right policy depends on the interconnect
    and active working-set ratio, not on the workload alone.

We implement {\uvmname} on PyTorch~2.8 with CUDA~13
    and evaluate Hugging~Face OPT and Qwen inference
    on A100-PCIe, H100-PCIe, and GH200 NVLink-C2C.
When memory is not oversubscribed, {\uvmname} matches PyTorch's native
    \texttt{cudaMalloc}-based allocator within $20\%$,
    while already outperforming default UVM by $52\times$--$58\times$
    and HF KV-cache offloading by $30\times$--$35\times$.
Under low oversubscription (ratio $<1.10$), where the native allocator crashes with OOM,
    {\uvmname} remains $45\times$ faster than default UVM and $27\times$ faster than HF offloading.
Under high oversubscription {\uvmname} still avoids OOM,
    but narrows to $1.6\times$--$1.7\times$ over default UVM
    and falls behind HF offloading by $\sim\!3\times$---a gap we are actively closing
    via the thrashing-mitigation mechanisms described above.
